\ifdefined\gkver\else\def\gkver{student}\fi
\documentclass[\gkver]{gaokaozhenti}
\begin{document}

\chapter{2020年北京}

\begin{enumerate}
\item
%% number: 1
%% paperName: 2020年北京等级考第 1 题 3 分
%% typeId: 1
%% type: 单选题
%% chapter: 光学
%% point: 光的干涉
%% method: 无
%% score: 3
%% degree: 950
%% duplicateId: 0
%% body:
以下现象不属于干涉的是\xzanswer{C}

\fourchoices[answer=C]
{白光经过杨氏双缝得到彩色图样}
{白光照射肥皂膜呈现彩色图样}
{白光经过三棱镜得到彩色图样}
{白光照射水面油膜呈现彩色图样}

%% answer: C
%% memo:
\memoanswer{
A．根据光的干涉定义可知白光经过杨氏双缝得到彩色图样是杨氏双缝干涉，故 A 错误；

B．由于重力的作用，肥皂膜形成了上薄下厚的薄膜，光线通过薄膜时频率不变，干涉条纹的产生是由于光线在薄膜前后两表面反射形成的两列光波的叠加，白光照射肥皂膜呈现彩色图样是属于干涉现象，故 B 错误；

C．白光经过三棱镜得到彩色图样是光在折射时产生的色散现象，故 C 正确；

D．水面上的油膜呈现彩色是光的干涉现象，属于薄膜干涉，故 D 错误；

故选 C。
}

\item
%% number: 2
%% paperName: 2020年北京等级考第 2 题 3 分
%% typeId: 1
%% type: 单选题
%% chapter: 原子和原子核
%% point: 玻尔模型
%% method: 无
%% score: 3
%% degree: 850
%% duplicateId: 0
%% body:
氢原子能级示意如图。现有大量氢原子处于 $n=3$ 能级上，下列说法正确的是\xzanswer{C}

\onepicture[w=4.00cm]{\includesvg{figs/02}}

\fourchoices[answer=C]
{这些原子跃迁过程中最多可辐射出 2 种频率的光子}
{从 $n=3$ 能级跃迁到 $n=1$ 能级比跃迁到 $n=2$ 能级辐射的光子频率低}
{从 $n=3$ 能级跃迁到 $n=4$ 能级需吸收 $0.66\UeV$ 的能量}
{$n=3$ 能级的氢原子电离至少需要吸收 $13.6\UeV$ 的能量}

%% answer: C
%% memo:
\memoanswer{
A．大量氢原子处于 $n=3$ 能级，最多可辐射出 $C_{3}^{2}=3$ 种不同频率的光子，故 A 错误；

B．根据能级图可知，从 $n=3$ 能级跃迁到 $n=1$ 能级辐射的光子能量为 $h\nu_{1}=13.6\UeV-1.51\UeV$，从 $n=3$ 能级跃迁到 $n=2$ 能级辐射的光子能量为 $h\nu_{2}=3.4\UeV-1.51\UeV$，比较可知从 $n=3$ 能级跃迁到 $n=1$ 能级比跃迁到 $n=2$ 能级辐射的光子频率高，故 B 错误；

C．根据能级图可知，从 $n=3$ 能级跃迁到 $n=4$ 能级需要吸收的能量为 $E=1.51\UeV-0.85\UeV=0.66\UeV$，故 C 正确；

D．根据能级图可知，氢原子处于 $n=3$ 能级的能量为 $-1.51\UeV$，故要使其电离至少需要吸收 $1.51\UeV$ 的能量，故 D 错误；

故选 C。
}

\item
%% number: 3
%% paperName: 2020年北京等级考第 3 题 3 分
%% typeId: 1
%% type: 单选题
%% chapter: 光学
%% point: 光的电磁说
%% method: 无
%% score: 3
%% degree: 950
%% duplicateId: 0
%% body:
随着通信技术的更新换代，无线通信使用的电磁波频率更高，频率资源更丰富，在相同时间内能够传输的信息量更大。第 5 代移动通信技术（简称 5G）意味着更快的网速和更大的网络容载能力，“4G 改变生活，5G 改变社会”。与 4G 相比，5G 使用的电磁波\xzanswer{A}

\fourchoices[answer=A]
{光子能量更大}
{衍射更明显}
{传播速度更大}
{波长更长}

%% answer: A
%% memo:
\memoanswer{
A．因为 5G 使用的电磁波频率比 4G 高，根据 $E=h\nu$ 可知 5G 使用的电磁波比 4G 光子能量更大，故 A 正确；

B．发生明显衍射的条件是障碍物（或孔）的尺寸可以跟波长相比，甚至比波长还小；因 5G 使用的电磁波频率更高，即波长更短，故 5G 越不容易发生明显衍射，故 B 错误；

C．光在真空中的传播速度都是相同的；光在介质中的传播速度为 $v=\frac{c}{n}$，5G 的频率比 4G 高，而频率越大折射率越大，光在介质中的传播速度越小，故 C 错误；

D．因 5G 使用的电磁波频率更高，根据 $v=\frac{c}{\lambda}$ 可知波长更短，故 D 错误；

故选 A。
}

\item
%% number: 4
%% paperName: 2020年北京等级考第 4 题 3 分
%% typeId: 1
%% type: 单选题
%% chapter: 气体性质
%% point: 气体图象
%% method: 无
%% score: 3
%% degree: 850
%% duplicateId: 0
%% body:
如图所示，一定量的理想气体从状态 A 开始，经历两个过程，先后到达状态 B 和 C。有关 A、B 和 C 三个状态温度 $T_{A}$、$T_{B}$ 和 $T_{C}$ 的关系，正确的是\xzanswer{C}

\onepicture[w=6.00cm]{\includesvg{figs/04}}

\fourchoices[answer=C]
{$T_{A}=T_{B}$，$T_{B}=T_{C}$}
{$T_{A}<T_{B}$，$T_{B}<T_{C}$}
{$T_{A}=T_{C}$，$T_{B}>T_{C}$}
{$T_{A}=T_{C}$，$T_{B}<T_{C}$}

%% answer: C
%% memo:
\memoanswer{
由题图可知，状态 A 到状态 B 为等压过程，根据盖—吕萨克定律 $\frac{V_{A}}{T_{A}}=\frac{V_{B}}{T_{B}}$，因为 $V_{B}>V_{A}$，故 $T_{B}>T_{A}$；状态 B 到状态 C 为等容过程，有 $\frac{p_{B}}{T_{B}}=\frac{p_{C}}{T_{C}}$，因为 $p_{B}>p_{C}$，故 $T_{B}>T_{C}$；对状态 A 和 C，由理想气体状态方程 $\frac{pV}{T}=C$ 可知 $T_{A}=T_{C}$。综上分析可知 C 正确，ABD 错误；

故选 C。
}

\item
%% number: 5
%% paperName: 2020年北京等级考第 5 题 3 分
%% typeId: 1
%% type: 单选题
%% chapter: 曲线运动
%% point: 人造地球卫星
%% method: 无
%% score: 3
%% degree: 650
%% duplicateId: 0
%% body:
我国首次火星探测任务被命名为“天问一号”。已知火星质量约为地球质量的 10\%，半径约为地球半径的 50\%，下列说法正确的是\xzanswer{A}

\fourchoices[answer=A]
{火星探测器的发射速度应大于地球的第二宇宙速度}
{火星探测器的发射速度应介于地球的第一和第二宇宙速度之间}
{火星的第一宇宙速度大于地球的第一宇宙速度}
{火星表面的重力加速度大于地球表面的重力加速度}

%% answer: A
%% memo:
\memoanswer{
A．当发射速度大于第二宇宙速度时，探测器将脱离地球的引力在太阳系的范围内运动，火星在太阳系内，所以火星探测器的发射速度应大于第二宇宙速度，故 A 正确；

B．第二宇宙速度是探测器脱离地球的引力到太阳系中的临界条件，当发射速度介于地球的第一和第二宇宙速度之间时，探测器将围绕地球运动，故 B 错误；

C．万有引力提供向心力，则有 $\frac{GMm}{R^{2}}=\frac{mv_{1}^{2}}{R}$，解得第一宇宙速度为 $v_{1}=\sqrt{\frac{GM}{R}}$，所以火星的第一宇宙速度为 $v_{\text{火}}=\sqrt{\frac{10\%}{50\%}}v_{\text{地}}=\frac{\sqrt{5}}{5}v_{\text{地}}$，所以火星的第一宇宙速度小于地球的第一宇宙速度，故 C 错误；

D．万有引力近似等于重力，则有 $\frac{GMm}{R^{2}}=mg$，解得星表面的重力加速度 $g_{\text{火}}=\frac{GM_{\text{火}}}{R_{\text{火}}^{2}}=\frac{10\%}{(50\%)^{2}}g_{\text{地}}=\frac{2}{5}g_{\text{地}}$，所以火星表面的重力加速度小于地球表面的重力加速度，故 D 错误；

故选 A。
}

\item
%% number: 6
%% paperName: 2020年北京等级考第 6 题 3 分
%% typeId: 1
%% type: 单选题
%% chapter: 机械波
%% point: 波与振动图象
%% method: 无
%% score: 3
%% degree: 550
%% duplicateId: 0
%% body:
一列简谐横波某时刻波形如图~\ref{2020北京06a}所示。由该时刻开始计时，质点 L 的振动情况如图~\ref{2020北京06b}所示。下列说法正确的是\xzanswer{B}

\twopicture[labelA=2020北京06a, labelB=2020北京06b, numstyle=chinese, wA=6.80cm, wB=6.80cm]
{figs/06a.png}
{figs/06b.png}

\fourchoices[answer=B]
{该横波沿 $x$ 轴负方向传播}
{质点 N 该时刻向 $y$ 轴负方向运动}
{质点 L 经半个周期将沿 $x$ 轴正方向移动}
{该时刻质点 K 与 M 的速度、加速度都相同}

%% answer: B
%% memo:
\memoanswer{
AB．由图可知乙质点 L 的振动情况，该时刻质点 L 向 $y$ 轴正方向振动。根据上下坡法或者平移法可知，该横波沿 $x$ 轴正方向传播，质点 N 该时刻向 $y$ 轴负方向运动，故 A 错误，故 B 正确；

C．质点 L 只在平衡位置附近 $y$ 轴方向上下振动，波传播时，质点不会沿 $x$ 轴正方向移动，故 C 错误；

D．该时刻质点 K 与 M 的速度为零，质点 K 加速度为 $-y$ 方向，质点 M 加速度为 $+y$ 方向，故 D 错误；

故选 B。
}

\item
%% number: 7
%% paperName: 2020年北京等级考第 7 题 3 分
%% typeId: 1
%% type: 单选题
%% chapter: 电场
%% point: 等势面
%% method: 无
%% score: 3
%% degree: 650
%% duplicateId: 0
%% body:
真空中某点电荷的等势面示意如图，图中相邻等势面间电势差相等。下列说法正确的是\xzanswer{B}

\onepicture[w=4.00cm]{\includesvg{figs/07}}

\fourchoices[answer=B]
{该点电荷一定为正电荷}
{P 点的场强一定比 Q 点的场强大}
{P 点电势一定比 Q 点电势低}
{正检验电荷在 P 点比在 Q 点的电势能大}

%% answer: B
%% memo:
\memoanswer{
A．正电荷和负电荷周围的等势面都为一组同心球壳，该点电荷不一定为正电荷，故 A 错误；

B．相邻等势面间电势差相等，P 点附近的等差等势面更加密集，故 P 点的场强一定比 Q 点的场强大，故 B 正确；

C．正电荷和负电荷周围的等势面都为一组同心球壳，若为正点电荷，则 P 点电势一定比 Q 点电势高，故 C 错误；

D．从等势面的情况无法判断该点电荷为正点电荷还是负点电荷，无法判断 P 点电势与 Q 点电势的高低，就无法判断正检验电荷在 P 点和在 Q 点的电势能的大小，故 D 错误；

故选 B。
}

\item
%% number: 8
%% paperName: 2020年北京等级考第 8 题 3 分
%% typeId: 1
%% type: 单选题
%% chapter: 磁场
%% point: 电流的磁场
%% method: 无
%% score: 3
%% degree: 650
%% duplicateId: 0
%% body:
如图所示，在带负电荷的橡胶圆盘附近悬挂一个小磁针。现驱动圆盘绕中心轴高速旋转，小磁针发生偏转。下列说法正确的是\xzanswer{B}

\onepicture[w=4.00cm]{figs/08.png}

\fourchoices[answer=B]
{偏转原因是圆盘周围存在电场}
{偏转原因是圆盘周围产生了磁场}
{仅改变圆盘的转动方向，偏转方向不变}
{仅改变圆盘所带电荷的电性，偏转方向不变}

%% answer: B
%% memo:
\memoanswer{
AB．小磁针发生偏转是因为带负电荷的橡胶圆盘高速旋转形成电流，而电流周围有磁场，磁场会对放入其中的小磁针有力的作用，故 A 错误，B 正确；

C．仅改变圆盘的转动方向，形成的电流的方向与初始相反，小磁针的偏转方向也与之前相反，故 C 错误；

D．仅改变圆盘所带电荷的电性，形成的电流的方向与初始相反，小磁针的偏转方向也与之前相反，故 D 错误；

故选 B。
}

\item
%% number: 9
%% paperName: 2020年北京等级考第 9 题 3 分
%% typeId: 1
%% type: 单选题
%% chapter: 交流电
%% point: 变压器
%% method: 无
%% score: 3
%% degree: 650
%% duplicateId: 0
%% body:
如图所示，理想变压器原线圈接在 $u=U_{m}\sin(\omega t+\varphi)$ 的交流电源上，副线圈接三个阻值相同的电阻 $R$，不计电表内电阻影响。闭合开关 S 后\xzanswer{A}

\onepicture[w=7.50cm]{figs/09.png}

\fourchoices[answer=A]
{电流表 A$_{2}$ 的示数减小}
{电压表 V$_{1}$ 的示数减小}
{电压表 V$_{2}$ 的示数不变}
{电流表 A$_{1}$ 的示数不变}

%% answer: A
%% memo:
\memoanswer{
开关 S 闭合时，副线圈总的电阻减小，由于变压器的匝数比和输入的电压都不变，所以输出的电压也不变，即 V$_{1}$ 示数不变，但因副线圈的总电阻减小，则副线圈的总电流增大，则原线圈的电流增大，故 A$_{1}$ 的示数变大；由于副线圈的电流增大，故串联在副线圈的电阻 $R$ 两端的电压增大，而副线圈的总电压不变，所以副线圈并联部分的电压减小，即 V$_{2}$ 的示数减小，故电流表 A$_{2}$ 的示数减小，故 A 正确，BCD 错误；

故选 A。
}

\item
%% number: 10
%% paperName: 2020年北京等级考第 10 题 3 分
%% typeId: 1
%% type: 单选题
%% chapter: 气体性质
%% point: 分子力
%% method: 无
%% score: 3
%% degree: 650
%% duplicateId: 0
%% body:
分子力 $F$ 随分子间距离 $r$ 的变化如图所示。将两分子从相距 $r=r_{2}$ 处释放，仅考虑这两个分子间的作用，下列说法正确的是\xzanswer{D}

\onepicture[w=5.50cm]{\includesvg{figs/10}}

\fourchoices[answer=D]
{从 $r=r_{2}$ 到 $r=r_{0}$ 分子间引力、斥力都在减小}
{从 $r=r_{2}$ 到 $r=r_{1}$ 分子力的大小先减小后增大}
{从 $r=r_{2}$ 到 $r=r_{0}$ 分子势能先减小后增大}
{从 $r=r_{2}$ 到 $r=r_{1}$ 分子动能先增大后减小}

%% answer: D
%% memo:
\memoanswer{
A．从 $r=r_{2}$ 到 $r=r_{0}$ 分子间引力、斥力都在增加，但斥力增加得更快，故 A 错误；

B．由图可知，在 $r=r_{0}$ 时分子力为零，故从 $r=r_{2}$ 到 $r=r_{1}$ 分子力的大小先增大后减小再增大，故 B 错误；

C．分子势能在 $r=r_{0}$ 时分子势能最小，故从 $r=r_{2}$ 到 $r=r_{0}$ 分子势能一直减小，故 C 错误；

D．从 $r=r_{2}$ 到 $r=r_{1}$ 分子势能先减小后增大，故分子动能先增大后减小，故 D 正确；

故选 D。
}

\item
%% number: 11
%% paperName: 2020年北京等级考第 11 题 3 分
%% typeId: 1
%% type: 单选题
%% chapter: 力物体的平衡
%% point: 摩擦力
%% method: 无
%% score: 3
%% degree: 650
%% duplicateId: 0
%% body:
某同学利用图~\ref{2020北京11a}所示装置研究摩擦力的变化情况。实验台上固定一个力传感器，传感器用棉线拉住物块，物块放置在粗糙的长木板上。水平向左拉木板，传感器记录的 $F-t$ 图像如图~\ref{2020北京11b}所示。下列说法正确的是\xzanswer{C}

\twopicture[labelA=2020北京11a, labelB=2020北京11b, numstyle=chinese, wA=6.00cm, wB=6.00cm]
{figs/11a.png}
{figs/11b.png}

\fourchoices[answer=C]
{实验中必须让木板保持匀速运动}
{图乙中曲线就是摩擦力随时间的变化曲线}
{最大静摩擦力与滑动摩擦力之比约为 10:7}
{只用图乙中数据可得出物块与木板间的动摩擦因数}

%% answer: C
%% memo:
\memoanswer{
AB．为了能研究摩擦力随时间的变化曲线，故物块一直要处于静止状态，则向左的摩擦力一直与向右轻绳的拉力平衡，图乙是向右轻绳的拉力随时间变化曲线，故图乙也可以反映摩擦力随时间变化的曲线，由图乙知向右轻绳的拉力先增大后减小，最后趋于不变，故物块先受静摩擦力作用后受滑动摩擦力作用，所以不需要让木板保持匀速运动，故 AB 错误；

C．由图可知，最大静摩擦力约为 $10\UN$，滑动摩擦力约为 $7\UN$，故最大静摩擦力与滑动摩擦力之比约为 10:7，故 C 正确；

D．根据 $F_{f}=\mu F_{N}$，$F_{N}=mg$，可知由于不知道物块的重力，故无法求物块与木板间的动摩擦因数，故 D 错误；

故选 C。
}

\item
%% number: 12
%% paperName: 2020年北京等级考第 12 题 3 分
%% typeId: 1
%% type: 单选题
%% chapter: 电路
%% point: 传感器
%% method: 无
%% score: 3
%% degree: 650
%% duplicateId: 0
%% body:
图~\ref{2020北京12a}表示某金属丝的电阻 $R$ 随摄氏温度 $t$ 变化的情况。把这段金属丝与电池、电流表串联起来（图~\ref{2020北京12b}），用这段金属丝做测温探头，把电流表的刻度改为相应的温度刻度，就得到了一个简易温度计。下列说法正确的是\xzanswer{B}

\twopicture[labelA=2020北京12a, labelB=2020北京12b, numstyle=chinese, wA=5.60cm, wB=5.00cm]
{figs/12a.png}
{figs/12b.png}

\fourchoices[answer=B]
{$t_{A}$ 应标在电流较大的刻度上，且温度与电流是线性关系}
{$t_{A}$ 应标在电流较大的刻度上，且温度与电流是非线性关系}
{$t_{B}$ 应标在电流较大的刻度上，且温度与电流是线性关系}
{$t_{B}$ 应标在电流较大的刻度上，且温度与电流是非线性关系}

%% answer: B
%% memo:
\memoanswer{
由甲图可知，$t_{A}$ 点对应的电阻阻值较小，由闭合电路欧姆定律知对应电路中的电流较大，故 $t_{A}$ 应标在电流较大的刻度上；而 $t_{B}$ 点对应的电阻阻值较大，由闭合电路欧姆定律知对应电路中的电流较小，故 $t_{B}$ 应标在电流较小的刻度上；由图甲得 $R=R_{0}+kt$，其中 $R_{0}$ 为图线的纵截距，由闭合电路欧姆定律得 $I=\frac{E}{R+R_{g}+r}$，联立解得 $t=\frac{E}{kI}-\frac{R_{0}+R_{g}+r}{k}$，可知 $t$ 与 $I$ 是非线性关系，故 B 正确，ACD 错误；

故选 B。
}

\item
%% number: 13
%% paperName: 2020年北京等级考第 13 题 3 分
%% typeId: 1
%% type: 单选题
%% chapter: 运动和力
%% point: 动量守恒定律
%% method: 无
%% score: 3
%% degree: 450
%% duplicateId: 0
%% body:
在同一竖直平面内，3 个完全相同的小钢球（1 号、2 号、3 号）悬挂于同一高度；静止时小球恰能接触且悬线平行，如图所示。在下列实验中，悬线始终保持绷紧状态，碰撞均为对心正碰。以下分析正确的是\xzanswer{D}

\onepicture[w=3.20cm]{\includesvg{figs/13}}

\fourchoices[answer=D]
{将 1 号移至高度 $h$ 释放，碰撞后，观察到 2 号静止、3 号摆至高度 $h$。若 2 号换成质量不同的小钢球，重复上述实验，3 号仍能摆至高度 $h$}
{将 1、2 号一起移至高度 $h$ 释放，碰撞后，观察到 1 号静止，2、3 号一起摆至高度 $h$，释放后整个过程机械能和动量都守恒}
{将右侧涂胶的 1 号移至高度 $h$ 释放，1、2 号碰撞后粘在一起，根据机械能守恒，3 号仍能摆至高度 $h$}
{将 1 号和右侧涂胶的 2 号一起移至高度 $h$ 释放，碰撞后，2、3 号粘在一起向右运动，未能摆至高度 $h$，释放后整个过程机械能和动量都不守恒}

%% answer: D
%% memo:
\memoanswer{
A．1 号球与质量不同的 2 号球相碰撞后，1 号球速度不为零，则 2 号球获得的动能小于 1 号球撞 2 号球前瞬间的动能，所以 2 号球与 3 号球相碰撞后，3 号球获得的动能也小于 1 号球撞 2 号球前瞬间的动能，则 3 号不可能摆至高度 $h$，故 A 错误；

B．1、2 号球释放后，三小球之间的碰撞为弹性碰撞，且三小球组成的系统只有重力做功，所以系统的机械能守恒，但整个过程中，系统所受合外力不为零，所以系统动量不守恒，故 B 错误；

C．1、2 号碰撞后粘在一起，为完全非弹性碰撞，碰撞过程有机械能损失，所以 1、2 号球再与 3 号球相碰后，3 号球获得的动能不足以使其摆至高度 $h$，故 C 错误；

D．碰撞后，2、3 号粘在一起，为完全非弹性碰撞，碰撞过程有机械能损失，且整个过程中，系统所受合外力不为零，所以系统的机械能和动量都不守恒，故 D 正确；

故选 D。
}

\item
%% number: 14
%% paperName: 2020年北京等级考第 14 题 3 分
%% typeId: 1
%% type: 单选题
%% chapter: 运动和力
%% point: 牛顿定律的应用
%% method: 无
%% score: 3
%% degree: 450
%% duplicateId: 0
%% body:
在无风的环境，某人在高处释放静止的篮球，篮球竖直下落；如果先让篮球以一定的角速度绕过球心的水平轴转动（如图）再释放，则篮球在向下掉落的过程中偏离竖直方向做曲线运动。其原因是，转动的篮球在运动过程中除受重力外，还受到空气施加的阻力 $f_{1}$ 和偏转力 $f_{2}$。这两个力与篮球速度 $v$ 的关系大致为：$f_{1}=k_{1}v^{2}$，方向与篮球运动方向相反；$f_{2}=k_{2}v$，方向与篮球运动方向垂直。下列说法正确的是\xzanswer{C}

\onepicture[w=6.50cm]{figs/14.png}

\fourchoices[answer=C]
{$k_{1}$、$k_{2}$ 是与篮球转动角速度无关的常量}
{篮球可回到原高度且角速度与释放时的角速度相同}
{人站得足够高，落地前篮球有可能向上运动}
{释放条件合适，篮球有可能在空中持续一段水平直线运动}

%% answer: C
%% memo:
\memoanswer{
A．篮球未转动时，篮球竖直下落，没有受到偏转力 $f_{2}$ 的作用，而篮球转动时，将受到偏转力 $f_{2}$ 的作用，所以偏转力 $f_{2}=k_{2}v$ 中的 $k_{2}$ 与篮球转动角速度有关，故 A 错误；

B．空气阻力一直对篮球做负功，篮球的机械能将减小，篮球的角速度也将减小，所以篮球没有足够的能量回到原高度，故 B 错误；

C．篮球下落过程中，其受力情况如下图所示。

\onepicture[w=4.00cm]{\includesvg{figs/14_an}}

篮球下落过程中，由受力分析可知，随着速度不断增大，篮球受到 $f_{1}$ 和 $f_{2}$ 的合力沿竖直方向的分力可能比重力大，可使篮球竖直方向的分速度减小为零或变成竖直向上，所以篮球可能向上运动，故 C 正确；

D．如果篮球的速度变成水平方向，则空气阻力的作用会使篮球速度减小，则篮球受到的偏转力 $f_{2}$ 将变小，不能保持 $f_{2}$ 与重力持续等大反向，所以不可能在空中持续一段水平直线运动，故 D 错误；

故选 C。
}

\item
%% number: 15
%% paperName: 2020年北京等级考第 15 题 9 分
%% typeId: 4
%% type: 实验题
%% chapter: 运动和力
%% point: 验证牛顿第二定律
%% method: 无
%% score: 9
%% degree: 550
%% duplicateId: 0
%% body:
在“探究加速度与物体受力、物体质量的关系”实验中，做如下探究：

\begin{enumerate}
	\item
	为猜想加速度与质量的关系，可利用图~\ref{2020北京15a}所示装置进行对比实验。两小车放在水平板上，前端通过钩码牵引，后端各系一条细线，用板擦把两条细线按在桌上，使小车静止。抬起板擦，小车同时运动，一段时间后按下板擦，小车同时停下。	对比两小车的位移，可知加速度与质量大致成反比。关于实验条件，下列正确的是：\tkanswer[2.20cm]{B}（选填选项前的字母）。
	
	（A）小车质量相同，钩码质量不同
	
	（B）小车质量不同，钩码质量相同
	
	（C）小车质量不同，钩码质量不同
	\item
	某同学为了定量验证（1）中得到的初步关系，设计实验并得到小车加速度 $a$ 与质量 $M$ 的 7 组实验数据，如下表所示。在图~\ref{2020北京15b}所示的坐标纸上已经描好了 6 组数据点，请将余下的一组数据描在坐标纸上，并作出 $a-\frac{1}{M}$ 图像 \tkanswer[2.20cm]{见详解}。
	
	\begin{center}
	\begin{tabular}{|c|c|c|c|c|c|c|c|}
	\hline
	次数 & 1 & 2 & 3 & 4 & 5 & 6 & 7 \\
	\hline
	$a/(\Umsq)$ & 0.62 & 0.56 & 0.48 & 0.40 & 0.32 & 0.24 & 0.15 \\
	\hline
	$M/\Ukg$ & 0.25 & 0.29 & 0.33 & 0.40 & 0.50 & 0.71 & 1.00 \\
	\hline
	\end{tabular}
	\end{center}
	\item
	在探究加速度与力的关系实验之前，需要思考如何测“力”。请在图~\ref{2020北京15c}中画出小车受力的示意图。为了简化“力”的测量，下列说法正确的是：\tkanswer[2.20cm]{AD}（选填选项前的字母）。
	
	A．使小车沿倾角合适的斜面运动，小车受力可等效为只受绳的拉力
	
	B．若斜面倾角过大，小车所受合力将小于绳的拉力
	
	C．无论小车运动的加速度多大，砂和桶的重力都等于绳的拉力
	
	D．让小车的运动趋近于匀速运动，砂和桶的重力才近似等于绳的拉力
\end{enumerate}

\threepicture[labelA=2020北京15a, labelB=2020北京15b, labelC=2020北京15c, numstyle=arabic, align=right, wA=11.00cm, wB=5.50cm, wC=6.50cm]
{figs/15a.png}
{figs/15b.png}
{figs/15c.png}

%% answer:
\jdanswer{
\begin{enumerate}
	\item
	B

	\item
	所作 $a-\frac{1}{M}$ 图像见详解。

	\item
	AD

\end{enumerate}
}
%% memo:
\memoanswer{
（1）为了探究加速度与质量的关系，必须控制小车所受拉力相同，而让小车的质量不同，所以钩码质量相同，故 B 正确。

（2）数据描点和 $a-\frac{1}{M}$ 图像如图所示。

\onepicture[w=6.00cm]{figs/15_an1.png}

（3）小车的受力示意图如图所示。

\onepicture[w=5.50cm]{figs/15_an2.png}

A．使小车沿倾角合适的斜面运动，小车所受重力沿斜面的分力刚好等于小车所受的摩擦力，则小车受力可等效为只受绳的拉力，故 A 正确；

B．若斜面倾角过大，重力沿斜面的分力大于摩擦力，小车所受合力将大于绳的拉力，不利于简化“力”的测量，故 B 错误；

C．由牛顿第二定律可知，无论小车运动的加速度多大，砂和桶的重力都大于绳的拉力，故 C 错误；

D．让小车的运动趋近于匀速运动时，砂和桶的重力才近似等于绳的拉力，故 D 正确。

故选 AD。
}

\item
%% number: 16
%% paperName: 2020年北京等级考第 16 题 9 分
%% typeId: 4
%% type: 实验题
%% chapter: 电路
%% point: 测电动势与内阻
%% method: 无
%% score: 9
%% degree: 450
%% duplicateId: 0
%% body:
用图~\ref{2020北京16a}、图~\ref{2020北京16b}所示的两种方法测量某电源的电动势和内电阻（约为 $1\UO$）。其中 $R$ 为电阻箱，电流表的内电阻约为 $0.1\UO$，电压表的内电阻约为 $3\UkO$。

\begin{enumerate}
	\item
	利用图~\ref{2020北京16a}实验电路测电源的电动势 $E$ 和内电阻 $r$，所测量的实际是图~\ref{2020北京16c}中虚线框所示“等效电源”的电动势 $E'$ 和内电阻 $r'$。若电流表内电阻用 $R_{A}$ 表示，请你用 $E$、$r$ 和 $R_{A}$ 表示出 $E'$、$r'$，并简要说明理由。
	\item
	某同学利用图像分析甲、乙两种方法中由电表内电阻引起的实验误差。在图~\ref{2020北京16d}中，实线是根据实验数据（图甲：$U=IR$，图乙：$I=\frac{U}{R}$）描点作图得到的 $U-I$ 图像；虚线是该电源的路端电压 $U$ 随电流 $I$ 变化的 $U-I$ 图像（没有电表内电阻影响的理想情况）。
	
	在图~\ref{2020北京16d}中，对应图甲电路分析的 $U-I$ 图像是：\tkanswer[1.50cm]{C}；对应图乙电路分析的 $U-I$ 图像是：\tkanswer[1.50cm]{A}。
	\item
	综合上述分析，为了减小由电表内电阻引起的实验误差，本实验应选择图~\ref{2020北京16a}、图~\ref{2020北京16b}中的\tkanswer[1.50cm]{乙}（填“甲”或“乙”）。
\end{enumerate}

\fourpicture[labelA=2020北京16a, labelB=2020北京16b, labelC=2020北京16c, labelD=2020北京16d, numstyle=arabic, numA=甲, numB=乙, numC=2, numD=3, align=right, wA=5.00cm, wB=5.00cm, wC=6.00cm, wD=11.00cm]
{figs/16a.png}
{figs/16b.png}
{figs/16c.png}
{figs/16d.png}

%% answer:
\jdanswer{
\begin{enumerate}
	\item
	$E'=E$，$r'=r+R_{A}$。

	\item
	C，A

	\item
	乙

\end{enumerate}
}
%% memo:
\memoanswer{
（1）将电源和电流表视为等效电源，电源电动势是电源本身具有的属性，电流表不具有产生电动势的本领，所以等效电源的电动势仍然为 $E'=E$，而电流表的内阻和电源的内阻作为等效电源的内阻，即 $r'=r+R_{A}$。

（2）对甲图，考虑电表内阻时，根据闭合电路欧姆定律得 $E=U_{\text{路}}+Ir_{\text{内}}=U+I(r+R_{A})$，变形得 $U=-(r+R_{A})I+E$；直接通过实验获得数据，可得 $U=-rI+E$。图像与纵轴截距均为电源电动势 $E$，虚线对应的斜率大小为 $r$，实线对应的斜率大小为 $r+R_{A}$，所以对应图甲电路分析的 $U-I$ 图像是 C；

对乙图，考虑电表内阻时（即虚线对应的真实情况），根据闭合电路欧姆定律得 $E=U_{\text{路}}+Ir_{\text{内}}=U+(I+\frac{U}{R_{v}})r=U+Ir+U\frac{r}{R_{v}}$，变形得 $U=-\frac{R_{v}r}{R_{v}+r}I+\frac{R_{v}}{R_{v}+r}E$；直接通过实验获得数据，可得 $U=-rI+E$。虚线对应的斜率大小为 $r$，实线对应的斜率大小为 $\frac{R_{v}}{R_{v}+r}r<r$，虚线对应的纵轴截距为 $E$，实线对应的纵轴截距为 $\frac{R_{v}}{R_{v}+r}E<E$；两图线在 $U=0$ 时，对应的短路电流均为 $I_{\text{短}}=\frac{E}{r}$，所以对应图乙电路分析的 $U-I$ 图像是 A。

（3）图甲虽然测量的电源电动势准确，但电流表分压较为明显，所以内阻测量的误差很大；图乙虽然电动势和内阻测量均偏小，但是电压表内阻很大，分流不明显，所以电动势和内阻的测量误差较小，所以选择图乙可以减小由电表内电阻引起的实验误差。
}

\item
%% number: 17
%% paperName: 2020年北京等级考第 17 题 9 分
%% typeId: 6
%% type: 计算题
%% chapter: 曲线运动
%% point: 平抛运动
%% method: 无
%% score: 9
%% degree: 650
%% duplicateId: 0
%% body:
无人机在距离水平地面高度 $h$ 处，以速度 $v_{0}$ 水平匀速飞行并释放一包裹，不计空气阻力，重力加速度为 $g$。

\begin{enumerate}
	\item
	求包裹释放点到落地点的水平距离 $x$；
	\item
	求包裹落地时的速度大小 $v$；
	\item
	以释放点为坐标原点，初速度方向为 $x$ 轴方向，竖直向下为 $y$ 轴方向，建立平面直角坐标系，写出该包裹运动的轨迹方程。
\end{enumerate}

%% answer:
\jdanswer{
\begin{enumerate}
	\item
	$v_{0}\sqrt{\frac{2h}{g}}$

	\item
	$\sqrt{v_{0}^{2}+2gh}$

	\item
	$y=\frac{g}{2v_{0}^{2}}x^{2}$

\end{enumerate}
}
%% memo:
\memoanswer{
（1）包裹脱离无人机后做平抛运动，在竖直方向做自由落体运动，则 $h=\frac{1}{2}gt^{2}$，解得 $t=\sqrt{\frac{2h}{g}}$；水平方向上做匀速直线运动，所以水平距离为 $x=v_{0}t=v_{0}\sqrt{\frac{2h}{g}}$。

（2）包裹落地时，竖直方向速度为 $v_{y}=gt=g\sqrt{\frac{2h}{g}}$，落地时速度为 $v=\sqrt{v_{0}^{2}+v_{y}^{2}}=\sqrt{v_{0}^{2}+2gh}$。

（3）包裹做平抛运动，分解位移有 $x=v_{0}t'$，$y=\frac{1}{2}gt'^{2}$，两式消去时间得包裹的轨迹方程为 $y=\frac{g}{2v_{0}^{2}}x^{2}$。
}

\item
%% number: 18
%% paperName: 2020年北京等级考第 18 题 9 分
%% typeId: 6
%% type: 计算题
%% chapter: 电磁感应
%% point: 法拉第电磁感应定律
%% method: 无
%% score: 9
%% degree: 650
%% duplicateId: 0
%% body:
如图~\ref{2020北京18a}所示，$N=200$ 匝的线圈（图中只画了 2 匝），电阻 $r=2\UO$，其两端与一个 $R=48\UO$ 的电阻相连，线圈内有指向纸内方向的磁场。线圈中的磁通量按图~\ref{2020北京18b}所示规律变化。

\begin{enumerate}
	\item
	判断通过电阻 $R$ 的电流方向；
	\item
	求线圈产生的感应电动势 $E$；
	\item
	求电阻 $R$ 两端的电压 $U$。
\end{enumerate}

\twopicture[labelA=2020北京18a, labelB=2020北京18b, numstyle=chinese, align=right, wA=5.00cm, wB=5.50cm]
{figs/18a.png}
{figs/18b.png}

%% answer:
\jdanswer{
\begin{enumerate}
	\item
	$a\to b$

	\item
	$10\UV$

	\item
	$9.6\UV$

\end{enumerate}
}
%% memo:
\memoanswer{
（1）根据图像可知，线圈中垂直于纸面向里的磁场增大，为了阻碍线圈中磁通量的增大，根据楞次定律可知线圈中感应电流产生的磁场垂直于纸面向外，根据安培定则可知线圈中的感应电流为逆时针方向，所以通过电阻 $R$ 的电流方向为 $a\to b$。

（2）根据法拉第电磁感应定律有 $E=N\frac{\Delta\Phi}{\Delta t}=200\times\frac{0.015-0.010}{0.10}\UV=10\UV$。

（3）电阻 $R$ 两端的电压为路端电压，根据分压规律可知 $U=\frac{R}{R+r}E=\frac{48}{48+2}\times10\UV=9.6\UV$。
}

\item
%% number: 19
%% paperName: 2020年北京等级考第 19 题 10 分
%% typeId: 1
%% type: 单选题
%% chapter: 电场
%% point: 带电粒子的直线运动
%% method: 无
%% score: 10
%% degree: 450
%% duplicateId: 0
%% body:
如图~\ref{2020北京19a}所示，真空中有一长直细金属导线 MN，与导线同轴放置一半径为 $R$ 的金属圆柱面。假设导线沿径向均匀射出速率相同的电子，已知电子质量为 $m$，电荷量为 $e$。不考虑出射电子间的相互作用。

\begin{enumerate}
	\item
	可以用以下两种实验方案测量出射电子的初速度：
	
	（A）在柱面和导线之间，只加恒定电压；
	
	（B）在柱面内，只加与 MN 平行的匀强磁场。
	
	当电压为 $U_{0}$ 或磁感应强度为 $B_{0}$ 时，刚好没有电子到达柱面。分别计算出射电子的初速度 $v_{0}$。
	\item
	撤去柱面，沿柱面原位置放置一个弧长为 $a$、长度为 $b$ 的金属片，如图~\ref{2020北京19b}所示。在该金属片上检测到出射电子形成的电流为 $I$，电子流对该金属片的压强为 $p$。求单位长度导线单位时间内出射电子的总动能。
\end{enumerate}

\twopicture[labelA=2020北京19a, labelB=2020北京19b, numstyle=chinese, align=right, wA=3.60cm, wB=3.60cm]
{figs/19a.png}
{figs/19b.png}

%% answer: BI
\jdanswer{
\begin{enumerate}
	\item
	$v_{0}=\sqrt{\frac{2eU_{0}}{m}}$（方案 A）或 $v_{0}=\frac{eB_{0}R}{2m}$（方案 B）

	\item
	$\frac{\pi Rp^{2}eab}{mI}$

\end{enumerate}
}
%% memo:
\memoanswer{
（1）（A）在柱面和导线之间，只加恒定电压 $U_{0}$，粒子刚好没有电子到达柱面，此时速度为零，根据动能定理有 $-eU_{0}=-\frac{1}{2}mv_{0}^{2}$，解得 $v_{0}=\sqrt{\frac{2eU_{0}}{m}}$。

（B）在柱面内，只加与 MN 平行的匀强磁场，磁感应强度为 $B_{0}$ 时，刚好没有电子到达柱面，设粒子的偏转半径为 $r$，根据几何关系有 $2r=R$，根据洛伦兹力提供向心力，则有 $B_{0}ev_{0}=m\frac{v_{0}^{2}}{r}$，解得 $v_{0}=\frac{eB_{0}R}{2m}$。

（2）设单位长度导线在单位时间内出射电子数为 $n$，金属片上电子形成的电流为 $I$，由电流的定义式得 $I=\frac{Q}{\Delta t}=\frac{nabe\cdot\Delta t}{2\pi R\cdot\Delta t}=\frac{nabe}{2\pi R}$，解得 $n=\frac{2\pi RI}{abe}$；电子流对金属片的压力大小为 $F=pS=pab$，对射到金属片上的电子，由动量定理得 $-pab\cdot\Delta t=-\frac{nab\cdot\Delta t}{2\pi R}\cdot mv$，解得 $v=\frac{peab}{mI}$，故单位长度导线在单位时间内出射电子的总动能为 $E_{k}=\frac{1}{2}mv^{2}\cdot n=\frac{\pi Rp^{2}eab}{mI}$。

}

\item
%% number: 20
%% paperName: 2020年北京等级考第 20 题 12 分
%% typeId: 6
%% type: 计算题
%% chapter: 电磁感应
%% point: 电磁感应中的变加速
%% method: 无
%% score: 12
%% degree: 450
%% duplicateId: 0
%% body:
某试验列车按照设定的直线运动模式，利用计算机控制制动装置，实现安全准确地进站停车。制动装置包括电气制动和机械制动两部分。图~\ref{2020北京20a}所示为该列车在进站停车过程中设定的加速度大小 $a_{\text{车}}$ 随速度 $v$ 的变化曲线。

（注意：解题过程中需要用到、但题目没有给出的物理量，要在解题时做必要的说明）

\begin{enumerate}
	\item
	求列车速度从 $20\Ums$ 降至 $3\Ums$ 经过的时间 $t$ 及行进的距离 $x$。
	\item
	有关列车电气制动，可以借助图~\ref{2020北京20b}模型来理解。图中水平平行金属导轨处于竖直方向的匀强磁场中，回路中的电阻阻值为 $R$，不计金属棒 MN 及导轨的电阻。MN 沿导轨向右运动的过程，对应列车的电气制动过程，可假设 MN 棒运动的速度与列车的速度、棒的加速度与列车电气制动产生的加速度成正比。列车开始制动时，其速度和电气制动产生的加速度大小对应图~\ref{2020北京20a}中的 P 点。论证电气制动产生的加速度大小随列车速度变化的关系，并在图~\ref{2020北京20a}中画出图线。
	\item
	制动过程中，除机械制动和电气制动外，列车还会受到随车速减小而减小的空气阻力。分析说明列车从 $100\Ums$ 减到 $3\Ums$ 的过程中，在哪个速度附近所需机械制动最强？
\end{enumerate}

\twopicture[labelA=2020北京20a, labelB=2020北京20b, numstyle=arabic, align=right, wA=6.50cm, wB=7.00cm]
{figs/20a.png}
{figs/20b.png}

%% answer:
\jdanswer{
\begin{enumerate}
	\item
	$t\approx24.3\Us$，$x\approx279.3\Um$

	\item
	列车电气制动产生的加速度与列车的速度成正比，为过 P 点的正比例函数，论证过程见解析，画出的图线见详解。

	\item
	$3\Ums$

\end{enumerate}
}
%% memo:
\memoanswer{
（1）由题图~\ref{2020北京20a}可知，列车速度从 $v_{1}=20\Ums$ 降至 $v_{2}=3\Ums$ 的过程中，加速度大小 $a_{1}=0.7\Umsq$ 保持不变，由匀变速直线运动规律得 $v_{2}=v_{1}-a_{1}t$，$-2a_{1}x=v_{2}^{2}-v_{1}^{2}$，由加速度的定义式 $a=\frac{\Delta v}{\Delta t}$ 得 $t=\frac{\Delta v}{a}=\frac{20-3}{0.7}\Us=\frac{170}{7}\Us\approx24.3\Us$，由速度位移公式 $v^{2}-v_{0}^{2}=-2ax$ 得 $x=\frac{v^{2}-v_{0}^{2}}{-2a}=\frac{3^{2}-20^{2}}{-2\times0.7}\Um=279.3\Um$。

（2）由 MN 沿导轨向右运动切割磁感线产生感应电动势 $E=Blv$，回路中感应电流 $I=\frac{E}{R}$，MN 受到的安培力 $F=BIL$，加速度为 $a=\frac{F}{m}$，结合上面几式得 $a=\frac{B^{2}l^{2}v}{mR}$，所以棒的加速度与棒的速度为正比例函数。又因为列车的电气制动过程，可假设 MN 棒运动的速度与列车的速度、棒的加速度与列车电气制动产生的加速度成正比，所以列车电气制动产生的加速度与列车的速度成正比，为过 P 点的正比例函数。画出的图线如下图所示。

\onepicture[w=6.00cm]{figs/20_an.png}

（3）由（2）可知，列车速度越小，电气制动的加速度越小。由题设可知列车还会受到随车速减小而减小的空气阻力，所以电气制动和空气阻力产生的加速度都随速度的减小而减小。由题图~\ref{2020北京20a}中，列车速度从 $20\Ums$ 降至 $3\Ums$ 的过程中加速度大小 $a_{\text{车}}$ 随速度 $v$ 减小而增大，所以列车速度从 $20\Ums$ 降至 $3\Ums$ 的过程中所需的机械制动逐渐变强，所以列车速度为 $3\Ums$ 附近所需机械制动最强。
}

\end{enumerate}

\end{document}
